\documentclass[10pt,a4paper]{amsart}
\usepackage[margin=19mm]{geometry}
\usepackage{amsmath,amssymb,mathtools}
\usepackage[scr=rsfs,cal=euler]{mathalfa}
\usepackage{xfrac}
\usepackage[unicode,hidelinks,bookmarks=false]{hyperref}
\newcommand{\ie}{i.e.\ }
\newcommand{\cf}{cf.\ }
\newcommand{\resp}{respectively\ }
\newcommand{\Subsection}{\subsection}
\providecommand{\nobreakdash}{\nobreak\hbox{-}}
\setlength{\emergencystretch}{2em}
\multlinegap0pt
\usepackage{amssymb}
\usepackage[shortlabels]{enumitem}
\usepackage{tikz}
\newtheorem{lemma}{Lemma}
\newtheorem{theorem}{Theorem}
\newtheorem{definition}{Definition}
\title{On some inequalities for weighted products and ratios of the sine function}
\author{Augustine L. Mahu, Beno\^it F. Sehba, Cecilia D. Williams}
\date{Source: arXiv:2607.26092v1 (CC BY 4.0)}
\begin{document}
\maketitle

\noindent\emph{Source and licence.} On some inequalities for weighted products and ratios of the sine function. Version arXiv:2607.26092v1 (CC BY 4.0), distributed by arXiv under the Creative Commons Attribution 4.0 International licence, http://creativecommons.org/licenses/by/4.0/, as recorded in the arXiv licence metadata for this exact version and shown on the abstract page https://arxiv.org/abs/2607.26092v1. Full licence notice: \texttt{licenses/arxiv-2607.26092-CC-BY-4.0.txt}.\par\smallskip

\noindent\textit{Editorial scope.} Counted unit: the article's theorem thm:lower-main (source0.tex lines 290--299). The article's complete original proof of it is delivered with it (source0.tex lines 301--335, one \texttt{proof} environment of 35 lines). No mathematics is written, repaired or re-derived: the statement and the proof are the article's own lines, in the article's own notation, and no retained line is substituted or re-flowed.  Retained uncounted as the source-local context the proof needs: lines 125--133; lines 219--225; lines 270--286.  Nothing was omitted from any retained span, so the package carries no omission record.  No retained line contains a citation, so the wrapper carries no bibliography and no bibliography entry.  No retained line contains a figure environment, an image-inclusion command or drawing code, so no image is delivered and the package declares no asset dependency.  Editorial formatting only, in addition: an AMS \texttt{amsart} preamble that restates the article's own declarations and macros used by the retained lines, namely the environments \texttt{lemma}, \texttt{theorem}, \texttt{definition} on separate counters, together with the standard packages those lines need.  The retained environments print in this document as Lemma 1, Theorem 1, Definition 1 in order of appearance, whereas the article prints the same items with the numbering of its own full document.  The complete native source of the article as distributed in the arXiv e-print for this exact version is bundled unchanged as \texttt{sources/206241/source0.tex}.
\begin{lemma}\label{lem:reverseholder}
Let $f_1, \dots, f_m$ be positive measurable functions on a measure space $(X, \mu)$.
Let $p_1, \dots, p_m \in \mathbb{R} \setminus \{0\}$ satisfy $\displaystyle\sum_{i=1}^m \frac{1}{p_i} = 1$.
If at least one $p_i \in (0,1)$, then
\[
\int_X \prod_{i=1}^m f_i(x) \, d\mu(x) \;\ge\; \prod_{i=1}^m \left( \int_X f_i(x)^{p_i} \, d\mu(x) \right)^{1/p_i},
\]
where for $p_i < 0$, the term $f_i^{p_i} = 1 / f_i^{|p_i|}$, and all integrals are assumed finite.
\end{lemma}

\begin{theorem}\label{thm:upper-main}
Let $\alpha_1,\ldots,\alpha_n > 0$ and $x_1,\ldots,x_n \in (0,1)$. Let
$E \subseteq \mathcal N$ be any reflection set with $0 < M_E(\vec x) < 1$. Then
\begin{equation}\label{eq:upperbound}
\prod_{i=1}^{n} \sin^{\alpha_i}(\pi x_i) \;\leq\; \sin^{\alpha}(\pi M_E(\vec x)).
\end{equation}
\end{theorem}

\begin{definition}\label{def:lower-reflection}
A \emph{reflection pair} $(E,F)$ is a choice of subsets $E \subseteq \mathcal I$ and
$F\subseteq\mathcal J$ such that for each $i\in \mathcal I$ and $j\in\mathcal J$, the associated selected angles satisfy
\[
a_i = \begin{cases} x_i, & i\in\mathcal I\setminus E\\ 1-x_i, & i\in E\end{cases}
\qquad (i\in\mathcal I),
\]
\[
c_j = \begin{cases} x_j, & j\in\mathcal J\setminus F\\ 1-x_j, & j\in F\end{cases}
\qquad (j\in\mathcal J),
\]
and the associated weighted reflection mean is given by
\begin{equation}\label{eq:MEFdef}
M_{E,F}(\vec x) = \frac{\displaystyle\sum_{i\in\mathcal I}\alpha_i x_i - \sum_{i\in E}\alpha_i(2x_i-1)
- \sum_{j\in\mathcal J}\alpha_j x_j + \sum_{j\in F}\alpha_j(2x_j-1)}{\alpha}.
\end{equation}
\end{definition}

\begin{theorem}\label{thm:lower-main}
Let $\alpha_1,\ldots,\alpha_n,\alpha$ be as above, and let $(E,F)$ be any reflection pair
with $0<M_{E,F}(\vec x)<1$. Then
\begin{equation}\label{eq:genineq}
\sin^{\alpha}\!\left(\pi M_{E,F}(\vec x)\right)
\;\leq\;
\frac{\displaystyle\prod_{i\in\mathcal I} \sin^{\alpha_i}(\pi x_i)}
     {\displaystyle\prod_{j\in\mathcal J} \sin^{\alpha_j}(\pi x_j)}.
\end{equation}
\end{theorem}

\begin{proof}
Define exponents $p_i = \alpha/\alpha_i$ for $i\in\mathcal I$ and $q_j = -\alpha/\alpha_j$
for $j\in\mathcal J$. Then $\sum_{i\in\mathcal I} 1/p_i + \sum_{j\in\mathcal J} 1/q_j = 1$,
and $p_{i_0}\in(0,1)$.

Put $M = M_{E,F}(\vec x)$ and $N = 1-M$. Using the integral representation of the Beta
function, we obtain
\begin{align*}
B(M,N) &= \int_0^1 t^{M-1}(1-t)^{-M}\,dt \\
&= \int_0^1 \prod_{i\in\mathcal I} \bigl[ t^{a_i}(1-t)^{1-a_i} \bigr]^{1/p_i}
         \prod_{j\in\mathcal J} \bigl[ t^{c_j}(1-t)^{1-c_j} \bigr]^{1/q_j}
         \cdot t^{-1}(1-t)^{-1} \, dt,
\end{align*}
with $a_i,c_j$ as in Definition~\ref{def:lower-reflection}. As before, the exponent of $t$
sums to $M-1$ and the exponent of $(1-t)$ sums to $-M$, by~\eqref{eq:MEFdef}.

Applying Lemma~\ref{lem:reverseholder}, we obtain
\[
B(M,N) \ge \prod_{i\in\mathcal I} \left( \int_0^1 t^{a_i-1}(1-t)^{-a_i} \, dt \right)^{1/p_i}
         \prod_{j\in\mathcal J} \left( \int_0^1 t^{c_j-1}(1-t)^{-c_j} \, dt \right)^{1/q_j}.
\]
As in the proof of Theorem~\ref{thm:upper-main}, every one of these integrals equals
$\pi/\sin(\pi x_i)$ (resp. $\pi/\sin(\pi x_j)$), regardless of the reflection choice. Hence,
\[
B(M,N) \ge \prod_{i\in\mathcal I} \left( \frac{\pi}{\sin(\pi x_i)} \right)^{\alpha_i/\alpha}
         \prod_{j\in\mathcal J} \left( \frac{\pi}{\sin(\pi x_j)} \right)^{-\alpha_j/\alpha}.
\]
Raising both sides to the power $\alpha>0$ and using $B(M,N)=\pi/\sin(\pi M)$ gives
\[
\left(\frac{\pi}{\sin(\pi M)}\right)^{\alpha} \ge \prod_{i\in\mathcal I} \left( \frac{\pi}{\sin(\pi x_i)} \right)^{\alpha_i}
         \prod_{j\in\mathcal J} \left( \frac{\pi}{\sin(\pi x_j)} \right)^{-\alpha_j}.
\]
This reduces to~\eqref{eq:genineq} since $\displaystyle\prod_{i\in\mathcal I}\pi^{\alpha_i}\prod_{j\in\mathcal J}\pi^{-\alpha_j}=\pi^\alpha$.
The proof is complete.
\end{proof}

\end{document}
